\section{The Fourth Reverence}

\begin{quotex}
Now I set forth man as asking himself: What is that in me which enables me to sacrifice the inmost lures of my impulses and all wishes that proceed from my nature, to a law which promises me no advantage in return and no penalty if I transgress it; which indeed, the more sternly it commands and the less it offers in return, the more I reverence it? This question stirs our whole soul in amazed wonder at the greatness and sublimity of the inner faculty in man and the insolubility of the mystery which it conceals (for the answer: `it is freedom,' would be tautological, because it is freedom itself that creates the mystery). We can never tire of directing our attention to it and admiring in ourselves a power which yields to no power of nature ... Here is what Archimedes wanted, but did not find: a firm point on which reason could place its lever, and that without applying it to the present or to a future world, but merely to its inner idea of freedom (which immovable moral law provides as a sure foundation) in order by its principles to set in motion the human will, even in opposition to all nature.
\flright{\textsc{Immanuel Kant}}
\end{quotex}


\textbf{Johann Wolfgang Goethe}'s theory of the three reverences includes reverence

\begin{enumerate}


\item for that which is above us

\item for that which is around us

\item for that which is below us
\end{enumerate}


From these, arise the three kinds of genuine religion.

These reverences are summarized by \textbf{Thomas Carlyle}:

\begin{quotex} 
Goethe practically distinguishes the kinds of religion that are, or have been, in the world; and says that for men there are three reverences. The boys are all trained to go through certain gesticulations; to lay their hands on their breasts and look up to heaven, in sign of the first reverence; other forms for the other two: so they give their three reverences.

\begin{enumerate}


\item The first and simplest is that of reverence for what is above us. It is the soul of all the Pagan religion; there is nothing better in the antique man than that.

\item Then there is reverence for what is around us,---reverence for our equals, to which he attributes an immense power in the culture of man.

\item The third is reverence for what is beneath us; to learn to recognise in pain, in sorrow and contradiction, even in those things, odious to flesh and blood, what divine meanings are in them; to learn that there lies in these also, and more than in any of the preceding, a priceless blessing. And he defines that as being the soul of the Christian religion,---the highest of all religions; `a height,' as Goethe says, `a height to which mankind was fated and enabled to attain; and from which, having once attained it, they can never retrograde.' Man cannot quite lose that (Goethe thinks), or permanently descend below it again; but always, even in the most degraded, sunken and unbelieving times, he calculates there will be found some few souls who will recognise what this highest of the religions meant; and that, the world having once received it, there is no fear of its ever wholly disappearing.
\end{enumerate}
\end{quotex}


But Goethe added a \textbf{fourth reverence}: true religion arises from the ``highest reverence'', that is, reverence for Self; it is only when he has reached this stage that man attains the highest pinnacle that he is capable of attaining.

But let us be careful to distinguish between this reverence for Self and attachment to ego. The ego is part of nature, snarled in the web of cause and effect. As Kant realized quoted above , this Self that Goethe speaks of is transcendent to nature, and only there can it be free.

Some New Age speakers speak about getting in touch with our ``Source'' and how all these wonderful things will happen. Fine as far as it goes, but it stops short, since it serves only to feed the ego. This ego can no more get in touch with its source than the Maid of the Mist\footnote{\url{https://www.niagarafallscrowneplazahotel.com/blog/2015/01/real-maid-mist-niagaras-famous-legend/}} climb back up to the top of Niagara Falls.

The talk feeds the desire of the ego to be the Infinite, a task it can only fail in -- it is the Self that is Infinite. The human being is a mix of Being and Nonbeing, act and potency, the actual and the possible. As such, he has limitations and opportunities, both of which he tends to underestimate.

As act, the human is never some abstract "human being", but is always someone in particular --- of some gender, culture, language, ethnicity, race --- it is this particularity that gives us our dharma. And the real opportunity of the human state, not available to other creatures in our world, is not to take up some new sport or manifest a parking space, but rather to transcend totally the human state, finally to be free.

When you seek to manifest your destiny, are you stirred in your ``whole soul in amazed wonder at the greatness and sublimity of the inner faculty in man and the insolubility of the mystery which it conceals''? Do you seek to accumulate things or seek the ideal good?

Nurture this feeling of awe and reverence for the ``Source'' --- don't try to get in touch with it. Rather, get out of its way, and let its light manifest in the world.


\flrightit{Posted on 2008-06-02 by Cologero}

